\documentclass[11pt]{amsart}
\usepackage[margin=1.1in]{geometry}
\usepackage{amssymb,mathtools}
\usepackage{booktabs,longtable}
\usepackage[hidelinks]{hyperref}
\usepackage{xurl}

\newtheorem{theorem}{Theorem}
\newtheorem{proposition}[theorem]{Proposition}
\theoremstyle{remark}
\newtheorem{remark}[theorem]{Remark}

\newcommand{\Q}{\mathbb{Q}}
\newcommand{\C}{\mathbb{C}}
\newcommand{\R}{\mathbb{R}}

\title{Planar sets with a fixed number of distances:\\ how the table was computed}
\author{Alon Amit}
\author{Claude (Anthropic)}
\date{November 2026}

\begin{document}

\begin{abstract}
The applet \emph{Planar sets with a fixed number of distances}
(\url{https://affinemess.com/#applets/fixed-distances}) tabulates, for $n\le16$ points and $d\le7$
distances, the planar $n$-point sets that determine exactly $d$ distinct distances, up to similarity, sorted
by the dimension of the family they belong to. This note records how each cell was obtained: which cells
come from the literature and how we re-verified them, which cells we computed for the first time and by what
method, what each computation assumes, and which files reproduce it. The classification of the sets with
fourteen, fifteen and sixteen points and seven distances is proved in two companion write-ups and only
summarised here.
\end{abstract}

\maketitle

\section{The table}\label{sec:table}

A configuration of $n$ labelled points modulo similarity has $2n-4$ degrees of freedom. Its \emph{distance
pattern} is the colouring of the edges of $K_n$ in which two pairs get the same colour exactly when they are
at the same distance; patterns are considered up to permutations of the points and of the colours. The
configurations with a given pattern form a semi-algebraic set; a \emph{family} is a realisable pattern
together with (a component of) this set, and its dimension is the dimension of the set modulo similarity.
Rigid configurations are the points of dimension $0$. Requiring exactly $d$ values among the $\binom n2$
distances imposes $\binom n2-d$ coincidences, so every component has dimension at least
$e(n,d)=2n-4-\binom n2+d$; symmetric configurations can have larger components, which is why some cells mix
rigid configurations with families.

A cell $(n,d)$ is \emph{solid} when its list is complete and \emph{dashed} when it shows only the
configurations we know (Section~\ref{sec:dashed}). Solid cells carry a mark that says where the proof of
completeness comes from:
\begin{itemize}
\item $\dagger$: a published classification, re-verified independently by an exact computation;
\item $\ddagger$: our exact computation only (we know of no published classification);
\item no mark: the three-point row, which needs no proof.
\end{itemize}
Before November 2026 several published cells were unmarked (the published lists had been reconstructed but
not recomputed). Section~\ref{sec:published} describes the computation that re-verified them; since then
every solid cell with at least four points carries $\dagger$ or $\ddagger$.

``Exact computation'' means: every candidate pattern was solved symbolically, so unrealisable patterns are
proved unrealisable, rigid configurations are isolated as real algebraic numbers, and positive-dimensional
components were decomposed and tested for real points. Section~\ref{sec:assume} lists what this relies on.

Table~\ref{tab:cells} lists the solid cells. The column ``families'' gives the number of families of each
dimension, the dimension in parentheses.

\begin{longtable}{@{}lrllp{0.42\textwidth}@{}}
\caption{The solid cells.}\label{tab:cells}\\
\toprule cell & rigid & families & mark & proof of completeness \\ \midrule
\endfirsthead
\toprule cell & rigid & families & mark & proof of completeness \\ \midrule
\endhead
$(3,1)$ & 1 & -- &  & trivial \\
$(3,2)$ & 0 & 1\,(1) &  & trivial \\
$(3,3)$ & 0 & 1\,(2) &  & trivial \\
$(4,2)$ & 6 & -- & $\dagger$ & \cite{ES}; \S\ref{sec:small} \\
$(4,3)$ & 0 & 9\,(1) & $\ddagger$ & \S\ref{sec:small} \\
$(4,4)$ & 0 & 7\,(2) & $\ddagger$ & \S\ref{sec:small} \\
$(4,5)$ & 0 & 2\,(3) & $\ddagger$ & \S\ref{sec:small} \\
$(4,6)$ & 0 & 1\,(4) & $\ddagger$ & \S\ref{sec:small} \\
$(5,2)$ & 1 & -- & $\dagger$ & \cite{ES}; \S\ref{sec:small} \\
$(5,3)$ & 34 & -- & $\dagger$ & \cite{Sh04}; \S\ref{sec:small} \\
$(5,4)$ & 468 & 25\,(1) & $\ddagger$ & \S\ref{sec:small} \\
$(5,5)$ & 0 & 402\,(1), 8\,(2) & $\ddagger$ & \S\ref{sec:small} \\
$(5,6)$ & 0 & 244\,(2), 1\,(3) & $\ddagger$ & \S\ref{sec:small} \\
$(5,7)$ & 0 & 76\,(3) & $\ddagger$ & \S\ref{sec:small} \\
$(6,3)$ & 9 & -- & $\dagger$ & \cite{Sh04}; \S\ref{sec:published} \\
$(6,4)$ & 149 & 4\,(1) & $\ddagger$ & \S\ref{sec:middle} \\
$(6,5)$ & 3062 & 38\,(1), 1\,(2) & $\ddagger$ & \S\ref{sec:middle} \\
$(6,6)$ & 44242 & 754\,(1), 3\,(2) & $\ddagger$ & \S\ref{sec:middle} \\
$(7,3)$ & 2 & -- & $\dagger$ & \cite{EF,Sh04}; \S\ref{sec:published} \\
$(7,4)$ & 42 & -- & $\dagger$ & \cite{LW}; \S\ref{sec:published} \\
$(7,5)$ & 508 & 9\,(1) & $\ddagger$ & \S\ref{sec:middle} \\
$(7,6)$ & 7161 & 51\,(1) & $\ddagger$ & \S\ref{sec:middle} \\
$(8,4)$ & 15 & -- & $\dagger$ & \cite{Sh08}; \S\ref{sec:published} \\
$(8,5)$ & 165 & 2\,(1) & $\ddagger$ & \S\ref{sec:middle} \\
$(8,6)$ & 1897 & 12\,(1) & $\ddagger$ & \S\ref{sec:middle} \\
$(9,4)$ & 4 & -- & $\dagger$ & \cite{EF}; \S\ref{sec:published} \\
$(9,5)$ & 57 & 1\,(1) & $\ddagger$ & \S\ref{sec:middle} \\
$(9,6)$ & 698 & 1\,(1) & $\ddagger$ & \S\ref{sec:middle} \\
$(10,5)$ & 19 & -- & $\ddagger$ & \S\ref{sec:ten}; \cite[Cor.~7]{AC} \\
$(10,6)$ & 238 & -- & $\ddagger$ & \S\ref{sec:ten}; \cite[Cor.~7]{AC} \\
$(11,5)$ & 4 & -- & $\dagger$ & \cite{W11}; \S\ref{sec:published} \\
$(11,6)$ & 64 & -- & $\dagger$ & \cite{AA}; \S\ref{sec:ten} \\
$(12,5)$ & 1 & -- & $\dagger$ & \cite{EF,Sh08}; \S\ref{sec:published} \\
$(12,6)$ & 14 & -- & $\ddagger$ & \S\ref{sec:ten} \\
$(13,6)$ & 3 & -- & $\dagger$ & \cite{W12,SO}; \S\ref{sec:published} \\
$(14,7)$ & 84 & -- & $\ddagger$ & \cite{AC14} \\
$(15,7)$ & 16 & -- & $\ddagger$ & \cite[Thm.~16]{AC} \\
$(16,7)$ & 2 & -- & $\ddagger$ & \cite[Thm.~1]{AC} \\
\bottomrule
\end{longtable}


\section{Candidate patterns}\label{sec:patterns}

Following Szöllősi and Östergård \cite{SO}, a pattern on $n$ points with at most six colours is a
\emph{candidate} if it is \emph{complex-realisable with distinct non-zero values}: with the squared
distances $s_{ij}$ written as one unknown per colour (the colour of one edge set to $1$), the ideal generated
by the $3\times3$ minors of the Gram matrix $C_{ij}=(s_{0i}+s_{0j}-s_{ij})/2$ together with the Rabinowitsch
equation $t\prod_{a<b}(s_a-s_b)\prod_c s_c=1$ is proper. The list $L_n$ is built hereditarily: the
candidates on $n$ points are the colourings of $K_n$ all of whose $(n-1)$-point sub-colourings lie in
$L_{n-1}$ and which pass the test. The pattern of every real $n$-point set with at most six distances lies in
$L_n$. Our re-implementation (nauty canonical forms, Gröbner bases modulo a prime) reproduces both rows of
Table~9 of \cite{SO} at every level:

\begin{center}
\begin{tabular}{@{}lrrrrrrrrrr@{}}
\toprule
$n$ & 4 & 5 & 6 & 7 & 8 & 9 & 10 & 11 & 12 & 13\\
\midrule
raw candidates & 25 & 1193 & 537\,230 & 14\,732 & 1565 & 635 & 228 & 62 & 15 & 3\\
$|L_n|$ & 24 & 922 & 20\,229 & 4667 & 1565 & 635 & 228 & 62 & 15 & 3\\
\bottomrule
\end{tabular}
\end{center}

For $n\le10$ the lists are \emph{certified} in \cite[Sections~2--3]{AC}: every pattern discarded along the
way carries a polynomial identity that proves the ideal is the unit ideal, and a short checker verifies all
$527\,338$ identities. The lists $L_7,\dots,L_{10}$ used below are identical, pattern for pattern, to the
certified ones. The lists $L_{11},L_{12},L_{13}$ come from the re-implementation only; their counts agree
with \cite{SO}.

\section{Four and five points}\label{sec:small}

For four and five points, and for every number of distances, each colouring type of $K_n$ (up to symmetry)
was solved in coordinates: $p_0=(0,0)$, $p_1=(x_1,0)$, the colour of the edge $p_0p_1$ equal to $1$, one
unknown per other colour, one equation per edge, and a Rabinowitsch variable for ``all values distinct and
non-zero''. The systems were handed to msolve \cite{msolve}: an empty Gröbner basis proves a type
unrealisable, and a zero-dimensional system yields its real solutions with certified isolating intervals.
Positive-dimensional systems were decomposed into minimal primes with Singular \cite{Singular}; each prime
was checked for degeneracy and probed for real points by random hyperplane slices.

\begin{itemize}
\item Four points: $5$ types with two distances give the six quadruples of \cite{ES}; the $9$ types with
  three distances, the $7$ with four, the $2$ with five and the single generic type are all realisable,
  each of the expected dimension, so the four-point cells consist of families only.
\item Five points: $17$ types with two colours give the regular pentagon; $124$ types with three colours
  give $35$ certified real roots, $34$ configurations up to similarity, Shinohara's list \cite{Sh04} (each
  matched to his figures 501--534 by its distance ratios; his table gives $2\sin4\pi/9$ for one ratio of
  figure 524 where the construction in his text gives $2\sin7\pi/18$). With four colours, $237$ of the
  $371$ types are realisable: $468$ rigid configurations and $25$ one-parameter families. With five, six
  and seven colours, $407$ of $443$, $244$ of $249$ and $76$ of $76$ types are realisable, and every
  component has dimension at least $1$, $2$, $3$ respectively.
\end{itemize}

\paragraph{Counting up to similarity.} The real roots of each type were grouped into orbits of the
colour-permuting automorphisms of the colouring (with rescaling) by interval arithmetic on the certified
isolating boxes; the real-root counts were confirmed by Sturm sequences on msolve's rational
parametrisations. (A first version reported $481$ rigid configurations for $(5,4)$; it compared distances
without rescaling and counted thirteen configurations twice.)

\paragraph{Hygiene.} Every unrealisability verdict was re-verified with Singular (\texttt{std} $=1$ with the
Rabinowitsch variable), and with msolve again for five or more colours. For $21$ of the $768$ types with five to seven colours the primary
decomposition did not finish; there the certificate is the dimension of the saturated ideal, certified
real points on random slices, and full rank of the Jacobian at those points. Separately, the $21$ components
on which random slices found no real point were settled one by one: twelve are empty (a generator is a
positive-definite quadratic), nine have real points on bounded pieces that the slices missed. No family
count changed.

\paragraph{Hidden rigid configurations.} A real point of a positive-dimensional type can be isolated in the
real locus without the dimension count seeing it. For five to seven distances we checked: wherever the
Jacobian of the ten equations has full rank the real locus is a manifold; at the finitely many real
rank-deficient points each component is smooth (after eliminating the linearly occurring variables, by
the minors of its Jacobian) or its singular locus has no real points. For four types with five distances
whose decomposition never finishes, the twenty candidate points were settled by a certified local test: a
smooth surface through the point contains the curve, and the intrinsic Hessian of a further generator is
indefinite there, so the point is a node of two real branches. None was found.

\paragraph{Structure of the solution sets.} For four and five points we also traced the one-parameter
families to their ends: rigid sets with fewer distances, collisions of two points, or infinity. For $(5,4)$
the families and their endpoints form a connected graph with $35$ vertices (the $34$ three-distance sets and
the pentagon), $108$ arcs, $77$ ends at a collision and $12$ at infinity; for $(5,5)$, $527$ vertices and
$5269$ arcs. The catalog cards list the arcs and the boundary strata of each family.

\section{Six to nine points with four to six distances}\label{sec:middle}

The nine cells $(6,4)$, $(6,5)$, $(6,6)$, $(7,5)$, $(7,6)$, $(8,5)$, $(8,6)$, $(9,5)$, $(9,6)$ were
classified from the candidate lists in the \emph{value space}. The unknowns are the $d-1$ free squared
distances, the equations the $3\times3$ minors of the Gram matrix (rank at most $2$), and the inequation
``values distinct and non-zero'' a Rabinowitsch variable. A full distance matrix determines a point set up
to isometry, so the configurations with a given pattern, modulo similarity, are the real solutions whose
Gram matrix is positive semidefinite of rank at most $2$ (rank $1$ means collinear).

For each pattern: msolve real root isolation when the system is zero-dimensional, followed by the
positive-semidefiniteness test (at $60$ digits) and reconstruction of the points; when it is
positive-dimensional, Singular's minimal primes, each prime tested for degeneracy, for real
positive-semidefinite points by random slices, and for real points of its singular locus (none occurred, so
every family is a smooth real curve or surface and no rigid configuration hides on it). Configurations are
counted up to similarity by an exact congruence test. About a dozen patterns of $(6,6)$ defeat
\texttt{minAssGTZ}; for them the ideal is split with \texttt{facstd} and each piece decomposed separately,
with Gröbner bases by \texttt{modStd}.

\begin{center}
\begin{tabular}{@{}lrrrrr@{}}
\toprule
cell & patterns & no real configuration & rigid & families & maximal rigid\\
\midrule
$(6,4)$ & 104 & 5 & 149 & 4 & 68\\
$(6,5)$ & 1509 & 116 & 3062 & 38 + 1 & 2084\\
$(6,6)$ & 18\,609 & 1351 & 44\,242 & 754 + 3 & 33\,304\\
$(7,5)$ & 362 & 21 & 508 & 9 & 184\\
$(7,6)$ & 4269 & 368 & 7161 & 51 & 3365\\
$(8,5)$ & 143 & 2 & 165 & 2 & 34\\
$(8,6)$ & 1408 & 73 & 1897 & 12 & 354\\
$(9,5)$ & 53 & 0 & 57 & 1 & 7\\
$(9,6)$ & 578 & 10 & 698 & 1 & 145\\
\bottomrule
\end{tabular}
\end{center}

\noindent ``Families'' counts one-parameter families, plus two-parameter ones after the plus sign;
``maximal'' counts the rigid sets contained in no set with one more point and the same number of distances.
The validation runs on $(5,3)$, $(5,4)$ and $(6,4)$ reproduced the counts of Section~\ref{sec:small} and of
an earlier coordinate-based computation of $(6,4)$.

\begin{remark}\label{rem:66}
In $(6,6)$, $19$ one-dimensional primes $P$ (of $911$ non-degenerate positive-dimensional ones, one in each of
$19$ patterns) showed no real positive-semidefinite point on $40$ random slices. Random slices certify real
points when they find them, but a miss is not a proof, so these $19$ were settled separately
(\texttt{cells/settle66.py}, a few minutes in all). Each pattern was decomposed again; the primes agree with the
first run, and in every pattern the primes other than these are exactly the families counted.

Let $C$ be the real points of $V(P)$, a closed subset of $\R^5$, and $K\subseteq C$ the points whose Gram
matrix is positive semidefinite: the real configurations, degenerate ones included. On $C$ the Gram matrix
has rank at most $2$, with eigenvalues $\lambda_a,\lambda_b,0,0,0$, so $K=\{\lambda_a\ge0,\lambda_b\ge0\}$ is
closed in $C$. Choose a rational point $A$ and four random rational combinations $g_1,\dots,g_4$ of the
generators of $P$, and let $W=P+\bigl(\det[\nabla g_1;\dots;\nabla g_4;x-A]\bigr)$; let
$R_1=P+(2\times2 \text{ minors of the Gram matrix})$. If $K$ meets a connected component $\Gamma$ of $C$, then
either $K\supseteq\Gamma$, and the point of $\Gamma$ nearest to $A$ (which exists since $\Gamma$ is closed) is a
critical point of the distance or a singular point of $V(P)$, hence a point of $W$ in $K$; or $K\cap\Gamma$ has a
boundary point in $\Gamma$, where $\lambda_a$ or $\lambda_b$ changes sign, so the Gram matrix has rank at most $1$
there: a point of $R_1$ in $K$. In every case $W$ and $R_1$ are zero-dimensional, and msolve isolates all their
real points. For $16$ of the $19$ primes none of these points is positive semidefinite, so $K=\emptyset$. For
the other $3$ some of them are, and a sign certificate finishes instead: the prime
contains $s_4+s_5$, $s_2+s_5$ or $s_3+2s_5$, so two squared distances of a configuration would have opposite
signs. Hence none of the $19$ curves carries a configuration, and the family count of $(6,6)$ is exact.
\end{remark}

\section{Ten to thirteen points with at most six distances}\label{sec:ten}

For $n=10,\dots,13$ every candidate pattern in $L_n$ has finitely many solutions (msolve; for $n=10$ and
$n=11$ also certified by Singular, $\dim\operatorname{std}(I)=0$ over $\Q$ for the coordinate system of
each pattern). The real solutions were computed by msolve on
Szöllősi--Östergård's own system (the Gram minors in the at most five free values, plus the Rabinowitsch
equation) with isolating intervals of width $2^{-400}$. For each real solution the Gram matrix was formed at
$100$ digits: it must be positive semidefinite of rank $2$ (two eigenvalues clearly positive, the others
about $10^{-100}$; none was indefinite), and the points come from its eigendecomposition; each reconstructed set is checked to have
exactly the pattern it came from. Up to similarity:

\begin{center}
\begin{tabular}{@{}lrrrr@{}}
\toprule
$n$ & patterns & five distances & six distances & total\\
\midrule
10 & 228 & 19 & 238 & 257\\
11 & 62 & 4 & 64 & 68\\
12 & 15 & 1 & 14 & 15\\
13 & 3 & 0 & 3 & 3\\
\bottomrule
\end{tabular}
\end{center}

These are the cells $(10,5)$, $(10,6)$, $(11,5)$, $(11,6)$, $(12,5)$, $(12,6)$, $(13,6)$. The ten-point
count was cross-checked by solving the much larger coordinate systems for $218$ of the $228$ patterns
(congruent realisations, class by class), and by checking that all $210$ ten-point sets with at most six
distances inside the twelvefold set $R_{12}+R_{12}$ (Section~\ref{sec:dashed}) and all ten-point subsets
of the $68$ eleven-point sets are among the $257$.

The row $n=10$ has since been proved independently in \cite[Corollary~7]{AC}: certified pattern list, exact
solution by SymPy Gröbner bases over number fields, and a published notebook that runs every step. The cells
$(11,5)$, $(12,5)$ and $(13,6)$ are published classifications \cite{W11,EF,Sh08,W12}; $(11,6)$ was first
classified in \cite[Lemma~6]{AA}, from the same $62$ candidate patterns, which is why it carries $\dagger$.
The cells $(10,5)$, $(10,6)$ and $(12,6)$ carry $\ddagger$.

\section{Fourteen to sixteen points with seven distances}\label{sec:top}

\begin{itemize}
\item $(16,7)$: $g(7)=16$, and the two Erdős--Fishburn lattice pieces are the only sixteen-point sets with
  at most seven distances \cite[Theorem~1]{AC}. The proof reduces every such set to one of the $257$
  ten-point sets above (or a regular polygon), with a new elementary proof that a $13$-cycle of diameters
  forces a regular $13$-gon, and extends each ten-point set by six points in exact arithmetic. A
  step-by-step notebook runs every verification \cite{ACnb}. That $g(7)=16$ was first proved in \cite{AA}.
\item $(15,7)$: exactly sixteen sets \cite[Theorem~16]{AC}, one of which (a lattice piece $F$) is not
  contained in any sixteen-point set.
\item $(14,7)$: exactly $84$ sets, all rigid \cite{AC14}.
\end{itemize}

An earlier, independent computation of the sixteen-point sets is described in the walkthrough
\cite{repro}: the same $257$ ten-point sets, extended by six points in $100$-digit arithmetic (circle
intersections and circumcentres as candidates, a depth-first search), give $826$ labelled sixteen-point sets
from $185$ cores, two up to similarity; the two patterns found are then solved exactly by msolve, each with a
unique realisation. It relied on the thirteen-cycle lemma of \cite{AA}, which \cite{AC} reproves.

\section{The published cells and their re-verification}\label{sec:published}

\begin{center}
\begin{tabular}{@{}llp{0.52\textwidth}@{}}
\toprule
cell & published & our check\\
\midrule
$(4,2)$, $(5,2)$ & \cite{ES} & Section~\ref{sec:small} ($5$ and $17$ types)\\
$(5,3)$ & \cite[Thm.~1]{Sh04} & Section~\ref{sec:small} ($124$ types)\\
$(6,3)$ & \cite[Thm.~2]{Sh04} & $7$ patterns of $L_6$ with three colours: $9$ sets\\
$(7,3)$ & \cite{EF,Sh04} & $2$ patterns of $L_7$: $2$ sets\\
$(7,4)$ & \cite[Thm.~8]{LW} & $34$ patterns of $L_7$: $42$ sets\\
$(8,4)$ & \cite[Thm.~1.2]{Sh08} & $14$ patterns of $L_8$: $15$ sets\\
$(9,4)$ & \cite{EF} & $4$ patterns of $L_9$: $4$ sets\\
$(11,5)$ & \cite[Thm.~13]{W11} & Section~\ref{sec:ten}: $4$ sets\\
$(12,5)$ & \cite{EF,Sh08} & Section~\ref{sec:ten}: $1$ set\\
$(13,6)$ & \cite{W12,SO} & Section~\ref{sec:ten}: $3$ sets\\
$(11,6)$ & \cite[Lemma~6]{AA} & Section~\ref{sec:ten}: $64$ sets\\
\bottomrule
\end{tabular}
\end{center}

The checks for $(6,3)$, $(7,3)$, $(7,4)$, $(8,4)$ and $(9,4)$ were added in November 2026
(\texttt{reverify\_published.py}). Each pattern of the certified list $L_n$ of \cite{AC} with exactly $d$
colours was solved with the routine \texttt{solve\_lengths} of \cite{AC} (a SymPy Gröbner basis of the Gram
minors with the Rabinowitsch equation, real roots in exact number fields); every pattern turned out to have
finitely many solutions, and the real Euclidean realisations were counted up to similarity. In every one of
the eleven cells the computed sets were then compared with the catalog: the counts agree, and the computed
sets match the catalog's sets one to one (compared by the multiset of the points' sorted distance profiles,
which separates all sets within each of these cells).

The maxima $g(2)=5$, $g(3)=7$, $g(4)=9$, $g(5)=12$ are due to Erdős and Fishburn \cite{EF}, who also
classified the maximum sets for $d\le4$; the twelve-point five-distance set is unique \cite{Sh08}, and
$g(6)=13$ \cite{W12}. The proofs of the top cells use the convex-polygon theorem of Altman and Fishburn
\cite{Al,Fi}.

\section{The dashed cells}\label{sec:dashed}

The dashed cells show the configurations we know; their counts are lower bounds, each configuration checked
numerically to have exactly the stated number of distances and the stated local dimension. The examples come
from subsets of regular polygons (with or without centre), pieces of the triangular and square lattices,
one-point extensions of the complete lists of smaller cells, a systematic search over configurations with a
cyclic or dihedral symmetry (orbit structures with up to five parameters and up to two extra
coincidences), and, for $8$--$12$ points with five distances, $9$--$13$ with six and $13$--$16$ with seven,
an exhaustive sweep of the subsets of the twelvefold set
$R_{12}+R_{12}=\{\zeta^a+\zeta^b:\zeta^{12}=1\}$ ($73$ points: the centre and points at angles that are
multiples of $15^\circ$ on six circles). This set contains pieces of the triangular lattice, the $12$-gon
with its centre, and twelvefold configurations that are neither, among them the three ten-point six-distance
patterns that no earlier family of examples realised. Configurations are counted up to similarity by an
exact congruence test (a vertex bijection matching all distances); the sweep contains many homometric pairs,
non-congruent sets with the same multiset of distances.

No dashed cell has a complete list: they are all open.

\section{What is exact and what is assumed}\label{sec:assume}

\begin{itemize}
\item \emph{Gröbner bases.} msolve computes over $\Q$ by multi-modular methods with probabilistic
  verification, like every current Gröbner engine; Singular's \texttt{std} and \texttt{minAssGTZ} work over
  $\Q$. The pattern lists for $n\le10$ do not depend on either: they are certified by explicit identities
  \cite{AC}.
\item \emph{Real roots.} msolve's isolating intervals are certified. The positive-semidefiniteness test is
  numerical, at $60$ to $100$ digits, with eigenvalues separated from zero by a wide margin in every case.
\item \emph{Real points on positive-dimensional components.} Found by random slices: a hit is a certified
  real point, a miss is not a proof. Every miss was then settled exactly: for five points by a
  positive-definite generator or a real point found by other means (Section~\ref{sec:small}), for $(6,6)$
  as in Remark~\ref{rem:66}; in the other cells there were no misses.
\item \emph{Similarity.} Counted by exact congruence tests or by interval arithmetic on certified roots.
\item \emph{Fourteen to sixteen points.} Exact arithmetic throughout; see \cite{AC,AC14}.
\item \emph{Cited theorems.} The classifications in Section~\ref{sec:published}, and Szöllősi--Östergård's
  counts, which our re-implementation reproduces and which serve as a check, not an input.
\end{itemize}

\section{Files}\label{sec:files}

Everything is in the folder \texttt{tools/fixed-distances} of the repository
\url{https://github.com/alonamit1729/affinemess-applets}; the companion proofs are in
\texttt{tools/16-7-proof}. The applet's data is \texttt{data/applet\_data.json}, assembled by
\texttt{build\_data2.py} and injected into the page by \texttt{build\_html.py}.

\begin{center}
\begin{tabular}{@{}lp{0.62\textwidth}@{}}
\toprule
step & scripts and results\\
\midrule
four and five points & \texttt{certify54.py}, \texttt{certify5d.py}, \texttt{certify\_nd.py}, \texttt{stubborn.py};
  \texttt{data/cert*.json}; hygiene: \texttt{reverify.py}, \texttt{dedup54.py}, \texttt{noreal.py},
  \texttt{noreal2.py}, \texttt{hidden.py}, \texttt{hidden\_final.py}, \texttt{nodetest.py}\\
structure & \texttt{strat.py}, \texttt{strat\_analyze.py}, \texttt{degen.py}, \texttt{degen\_analyze2.py}\\
candidate patterns & \texttt{g7fin/gen/gen.py}, \texttt{g7fin/gen/levels/}; certified lists:
  \texttt{tools/16-7-proof/certificates/}\\
six to nine points & \texttt{cells/cellcert.py}, \texttt{cells/summary.py}, \texttt{cells/export\_cells.py},
  \texttt{cells/settle66.py}\\
ten to thirteen points & \texttt{g7fin/pat\_rigid\_d.py}, \texttt{g7count/gram\_n.py},
  \texttt{g7count/sets10\_by\_d.json}, \texttt{g7count/sets1113\_by\_d.json}, \texttt{exact\_cells.py}\\
published cells & \texttt{reverify\_published.py}, \texttt{data/reverify\_published.json}\\
dashed cells & \texttt{symfam.py}, \texttt{extend.py}, \texttt{examples.py}, \texttt{lattice.py},
  \texttt{r12/sweep.py}\\
\bottomrule
\end{tabular}
\end{center}

\noindent Eleven reproduction-only files over 1\,MB are not in the published tree; \texttt{fetch\_data.py}
restores them from a pinned commit and checks their SHA-256. The walkthrough \cite{repro} gives the
commands and expected outputs for the ten- to thirteen-point computation and the earlier sixteen-point
computation.

\begin{thebibliography}{99}
\bibitem{AC} A. Amit and Claude (Anthropic), The largest planar seven-distance sets (November 2026),
  \url{https://affinemess.com/tools/16-7-proof/paper/seven-distances.pdf}.
\bibitem{ACnb} A. Amit and Claude (Anthropic), The $(16,7)$ proof: a step-by-step notebook,
  \url{https://affinemess.com/16-7-proof.html}; runnable in the browser at
  \url{https://alonamit1729.github.io/live-notebooks/16-7-proof.html}.
\bibitem{AC14} A. Amit and Claude (Anthropic), Fourteen points, seven distances (November 2026),
  \url{https://affinemess.com/tools/16-7-proof/fourteen/fourteen-points.pdf}.
\bibitem{repro} Reproducing the classification of ten- to sixteen-point sets (walkthrough),
  \url{https://affinemess.com/tools/fixed-distances/g7count/reproducing.html}.
\bibitem{AA} A. Amit and GPT-6 Astra, The maximum cardinality of a planar seven-distance set, preprint,
  Hexagon 2610.00534 (October 2026).
\bibitem{Al} E. Altman, On a problem of P. Erdős, \emph{Amer. Math. Monthly} \textbf{70} (1963), 148--157.
\bibitem{ES} S. J. Einhorn and I. J. Schoenberg, On Euclidean sets having only two distances between points
  I--II, \emph{Indag. Math.} \textbf{28} (1966).
\bibitem{EF} P. Erdős and P. Fishburn, Maximum planar sets that determine $k$ distances, \emph{Discrete
  Math.} \textbf{160} (1996), 115--125.
\bibitem{Fi} P. Fishburn, Convex polygons with few intervertex distances, \emph{Comput. Geom.} \textbf{5}
  (1995), 65--93.
\bibitem{LW} W. Lan and X. Wei, Classification of seven-point four-distance sets in the plane, \emph{Math.
  Notes} \textbf{93} (2013), 510--522.
\bibitem{msolve} J. Berthomieu, C. Eder and M. Safey El Din, msolve: a library for solving polynomial
  systems, \emph{ISSAC} (2021).
\bibitem{Sh04} M. Shinohara, Classification of three-distance sets in two dimensional Euclidean space,
  \emph{European J. Combin.} \textbf{25} (2004), 1039--1058.
\bibitem{Sh08} M. Shinohara, Uniqueness of maximum planar five-distance sets, \emph{Discrete Math.}
  \textbf{308} (2008), 3048--3055.
\bibitem{Singular} W. Decker, G.-M. Greuel, G. Pfister and H. Schönemann, \textsc{Singular} 4 --- a computer
  algebra system for polynomial computations.
\bibitem{SO} F. Szöllősi and P. R. J. Östergård, Constructions of maximum few-distance sets in Euclidean
  spaces, \emph{Electron. J. Combin.} \textbf{27}(1) (2020), P1.23.
\bibitem{W11} X. Wei, Classification of eleven-point five-distance sets in the plane, \emph{Ars Combin.}
  \textbf{102} (2011), 505--515.
\bibitem{W12} X. Wei, A proof of Erdős--Fishburn's conjecture for $g(6)=13$, \emph{Electron. J. Combin.}
  \textbf{19}(4) (2012), P38.
\end{thebibliography}

\end{document}
